\section{Fault observability and the one rank deficiency}
\label{sec:obs}

This section states the paper's structural core: the parity projector that survives the
rank-deficient, gauge-carrying $\Gmat$; the characterization of undetectable faults as
exactly $\rng(\Gmat)$; the one-rank-deficiency unification that ties the parameter-space
gauge to the measurement-space blind subspace; and the protection-gap consequence. The
classical parity result is \Cited{}; the gauge-survival, the unification, and the
protection gap are \Proven{}; the engine numbers that instantiate them are outputs of the
\Validated{} pipeline of Section~\ref{sec:anchor}.

\subsection{The weighted parity projector}
\label{sec:parity}

The weighted least-squares estimate of $\xvec$ from \eqref{eq:model} is
$\hat{\xvec}=(\Gmat^\top\Wmat\Gmat)^{+}\Gmat^\top\Wmat\bm z$, using the Moore--Penrose
pseudo-inverse because $\Gmat^\top\Wmat\Gmat$ is rank-deficient (its kernel is the gauge
\eqref{eq:gauge}). The residual is $\bm r=\Pperp\bm z$ with the \emph{weighted parity
projector}
\begin{equation}
  \Pperp \;=\; \bm{I} \;-\; \Gmat\,(\Gmat^\top\Wmat\Gmat)^{+}\,\Gmat^\top\Wmat
  \;\in\;\R^{84\times84}.
  \label{eq:pperp}
\end{equation}

\begin{proposition}[The parity projector is gauge-independent]
\label{prop:projector}
Let $\bm M:=\Wmat^{1/2}\Gmat$ and let $\Pi_{\bm M}$ be the orthogonal projector onto
$\rng(\bm M)$. Then $\Pperp=\Wmat^{-1/2}(\bm{I}-\Pi_{\bm M})\Wmat^{1/2}$ is idempotent and
$\Wmat$-self-adjoint (i.e.\ $\Wmat\Pperp=\Pperp{}^{\!\top}\Wmat$, \emph{not} ordinary
symmetric when $\Wmat\neq\bm{I}$), with $\Pperp\Gmat=\bm0$, $\ker\Pperp=\rng(\Gmat)$, and
$\rng(\Pperp)$ equal to the $\Wmat$-orthogonal complement of $\rng(\Gmat)$. Every one of
these objects depends only on $\rng(\Gmat)$ and $\Wmat$, hence is invariant under the
gauge: adding any $\bm n\in\Nsp(\Gmat)$ to the state leaves $\Pperp$ and $\bm r$ unchanged.
\end{proposition}

\begin{proof}
$\Gmat(\Gmat^\top\Wmat\Gmat)^{+}\Gmat^\top\Wmat$ is the $\Wmat$-orthogonal (oblique in the
standard inner product) projector onto $\rng(\Gmat)$: writing $\bm M=\Wmat^{1/2}\Gmat$, one
has $\Wmat^{1/2}\Gmat(\Gmat^\top\Wmat\Gmat)^{+}\Gmat^\top\Wmat^{1/2}=\bm M(\bm M^\top\bm
M)^{+}\bm M^\top=\Pi_{\bm M}$, the ordinary orthogonal projector onto $\rng(\bm M)$, whose
defining properties (idempotence, self-adjointness, range) are basis-free and survive the
rank deficiency because $(\bm M^\top\bm M)^{+}$ is the pseudo-inverse. Conjugating by
$\Wmat^{-1/2}$ gives \eqref{eq:pperp}, so $\Pperp=\Wmat^{-1/2}(\bm{I}-\Pi_{\bm M})\Wmat^{1/2}$
is idempotent, and $\Wmat\Pperp=\Wmat^{1/2}(\bm{I}-\Pi_{\bm M})\Wmat^{1/2}$ is symmetric,
i.e.\ $\Pperp$ is $\Wmat$-self-adjoint. Finally $\Pperp\Gmat=\Wmat^{-1/2}(\bm{I}-\Pi_{\bm
M})\bm M=\bm0$, so $\rng(\Gmat)\subseteq\ker\Pperp$; equality follows by dimension count,
$\dim\ker\Pperp=\rank\Gmat=\dim\rng(\Gmat)$. None of $\rng(\Gmat)$, $\Wmat$, or $\Pi_{\bm
M}$ references a basis of $\Nsp(\Gmat)$, so all are gauge-independent.
\end{proof}

Because $\mathbb E[\bm r]=\Pperp\,\mathbb E[\bm z]=\Pperp\bvec$ (the signal $\Gmat\xvec$ is
annihilated by $\Pperp$), the gauge does \emph{not} leak into detectability: the parity
residual sees a fault through $\Pperp$ alone.

\subsection{Undetectable faults are exactly \texorpdfstring{$\rng(\Gmat)$}{range(G)}}
\label{sec:undetectable}

\begin{theorem}[Undetectable-fault characterization, gauge-carrying $\Gmat$]
\label{thm:undetectable}
A fault $\bvec$ produces zero expected parity residual, $\mathbb E[\bm r]=\Pperp\bvec=\bm0$,
if and only if $\bvec\in\rng(\Gmat)$. The characterization holds for the rank-deficient,
gauge-carrying free-network $\Gmat$ through the pseudo-inverse projector \eqref{eq:pperp}.
Consequently, in a network with no redundancy ($n=\rank\Gmat$) one has $\rng(\Gmat)=\R^{n}$
and \emph{nothing} is detectable.
\end{theorem}

\begin{proof}
By Proposition~\ref{prop:projector}, $\ker\Pperp=\rng(\Gmat)$, so $\Pperp\bvec=\bm0\iff
\bvec\in\rng(\Gmat)$. The classical parity statement ``$\bvec$ undetectable
$\iff\bvec\in\rng(\Gmat)$'' is \Cited{} \citep{sturza1988,brown1992,parkinson1988}; the only
new content is that it survives a rank-deficient $\Gmat$, which Proposition~\ref{prop:projector}
supplies by using $(\Gmat^\top\Wmat\Gmat)^{+}$ in place of an inverse and by noting $\Pperp$
is gauge-independent. If $n=\rank\Gmat$ then $\rng(\Gmat)=\R^{n}$, so every $\bvec$ lies in
$\rng(\Gmat)$ and $\Pperp=\bm0$.
\end{proof}

The pipeline evaluates two witnesses on the real-DE440 network (both counts \Validated{},
Section~\ref{sec:anchor}). An \emph{in-range} fault $\bvec=\Gmat\,\bm e_{\dtau(\text{node
}1)}$---a shift of node~1's clock offset---has parity norm $\lVert\Pperp\bvec\rVert=9.581
\times10^{-12}$ and is \textbf{undetectable}. A \emph{generic} fault $\bvec=\bm e_0$ (a
single measurement bias) has $\lVert\Pperp\bvec\rVert=1.021$ and is \textbf{detectable}.

\subsection{Minimum detectable bias}
\label{sec:mdb}

For a unit fault direction $\bm c$ (a column of a fault signature), the Baarda minimum
detectable bias is the smallest fault magnitude a non-central parity test rejects at
prescribed false-alarm and detection powers \citep{baarda1968,teunissen2006}:
\begin{equation}
  \MDB(\bm c)\;=\;\sqrt{\frac{\lambda_0}{\bm c^\top\Wmat\Pperp\bm c}},
  \label{eq:mdb}
\end{equation}
where $\lambda_0$ is the non-centrality set by the false-alarm probability $P_{\mathrm{fa}}$
and the missed-detection probability $P_{\mathrm{md}}$. The non-centrality is the
\emph{single}-projector, $\Wmat$-weighted parity energy $\bm c^\top\Wmat\Pperp\bm c$; because
$\Pperp$ is $\Wmat$-self-adjoint and idempotent (Proposition~\ref{prop:projector}), it equals
$(\Pperp\bm c)^\top\Wmat(\Pperp\bm c)$ for \emph{all} $\Wmat$---so this correct form is the
weighted energy of the projected column. It is \emph{not} the double-projector form
$\bm c^\top\Pperp\Wmat\Pperp\bm c$, which differs whenever $\Wmat\neq\bm{I}$ and is the wrong
non-centrality; Appendix~\ref{app:mdb} states the identity and the counterexample precisely.
With $\lambda_0=17.075$ from $P_{\mathrm{fa}}=10^{-3}$, $P_{\mathrm{md}}=0.20$, the pipeline
reports (both \Validated{}) $\MDB(\bm e_0)=4.6592$ for the generic single-measurement fault,
and $\MDB=\infty$ for the in-range fault $\Gmat\,\bm e_{\dtau(\text{node }1)}$, whose
non-centrality $\bm c^\top\Wmat\Pperp\bm c=0$---the protection gap of
Proposition~\ref{prop:protection} read as a bias magnitude.

\subsection{The one rank deficiency, two faces}
\label{sec:unification}

\begin{proposition}[One rank deficiency, two faces]
\label{prop:unification}
The single rank deficiency of the self-referential Jacobian, $\rank\Gmat=32<40$, produces
\emph{both}
\begin{enumerate}
  \item[(a)] the parameter-space gauge $\Nsp(\Gmat)\subset\R^{40}$, of dimension
  $40-\rank\Gmat=8$ (the datum$\oplus$timescale gauge, P1$\,\oplus\,$P4); and
  \item[(b)] the measurement-space blind fault subspace $\rng(\Gmat)\subset\R^{84}$, of
  dimension $\rank\Gmat=32$ (the undetectable faults of Theorem~\ref{thm:undetectable}, with
  the P3 common-mode class a subset).
\end{enumerate}
The two are linked, not equal: for any undetectable fault $\bvec\in\rng(\Gmat)$ there is a
state shift $\Delta\xvec$ with $\Gmat\,\Delta\xvec=\bvec$, and $\Delta\xvec$ is determined
only \emph{modulo} $\Nsp(\Gmat)$. Thus an undetectable fault is precisely one absorbed into
a state error whose ambiguity is exactly the gauge.
\end{proposition}

\begin{proof}
The rank--nullity theorem applied to $\Gmat\colon\R^{40}\to\R^{84}$ gives
$\dim\Nsp(\Gmat)=40-\rank\Gmat=8$ and $\dim\rng(\Gmat)=\rank\Gmat=32$; both are functions of
the one integer $\rank\Gmat$, which is (a) and (b). For the link, $\bvec\in\rng(\Gmat)$ means
$\bvec=\Gmat\,\Delta\xvec$ has a solution by definition of the range; the solution set is the
coset $\Delta\xvec+\Nsp(\Gmat)$, since $\Gmat(\Delta\xvec+\bm n)=\Gmat\Delta\xvec$ for
$\bm n\in\Nsp(\Gmat)$ and only for such $\bm n$. Hence the induced state error is defined
modulo $\Nsp(\Gmat)$.
\end{proof}

This is the honest form of the program's unification: it is a statement about two subspaces
tied by one map, \emph{not} the (false) claim that the gauge and the blind subspace are the
same set. The pipeline reports the two dimensions directly---$\rank\Gmat=32$ and the parity
dimension $\operatorname{trace}(\Pperp)=84-32=52$ (both \Validated{}), against the eight
gauge directions $\Nsp(\Gmat)$ that Theorem P4 establishes (\Proven{}). The eight analytic
gauge generators are annihilated by $\Gmat$ to a worst relative residual of $2.702\times
10^{-17}$, confirming $\Nsp(\Gmat)\supseteq\operatorname{span}\{\text{3 transl},\text{3
rot},\dtau,\dalpha\}$ exactly. The count is spectrally unambiguous: the largest
gauge-direction eigenvalue of the (preconditioned) normal matrix is
$\lambda_{\text{null}}^{\max}=4.367\times10^{-15}$, the smallest observable eigenvalue is
$\lambda_{\text{obs}}^{\min}=5.455\times10^{-4}$, and the largest is
$\lambda^{\max}=8.597\times10^{1}$, a null-to-observable spectral gap of $1.249\times
10^{11}$.

\begin{remark}[Common-mode is a subset of the blind subspace]
\label{rem:commonmode}
The common-mode reference-frame error class of P3 \citep{baweja2026commonmode}---a shift
$\delta\bm y=\Gmat[\,-\Delta\bm s;\,\bm 0\,]$ that moves the frame realization---is of the
form $\Gmat\,(\cdot)$ and so lies in $\rng(\Gmat)$ by construction; it is a special case of
Theorem~\ref{thm:undetectable}, undetectable to any snapshot residual test. Paired with the
parameter-space gauge $\Nsp(\Gmat)$ (P1/P4), this is the two-faces statement of
Proposition~\ref{prop:unification}: the frame that a monitor cannot pin (parameter space)
and the frame-shaped fault it cannot see (measurement space) are the two images of the one
rank deficiency.
\end{remark}

\subsection{The protection gap}
\label{sec:protectiongap}

Theorem~\ref{thm:undetectable} has a protection-level consequence. For an undetectable fault
the residual is zero but the state estimate still moves; the ratio of the two is the Brown
slope, and for the $\rng(\Gmat)$ class it is infinite.

\begin{proposition}[The protection gap: infinite slope on $\rng(\Gmat)$]
\label{prop:protection}
Let $\Pi_{\mathrm{obs}}$ project the state error onto the observable (user-PVT) directions,
and define the weighted Brown slope of a fault as
\begin{equation}
  \operatorname{slope}(\bvec)
  \;=\;\frac{\lVert\Pi_{\mathrm{obs}}\,\Delta\hat{\xvec}\rVert}{\lVert\Pperp\bvec\rVert_{\Wmat}},
  \qquad
  \Delta\hat{\xvec}=(\Gmat^\top\Wmat\Gmat)^{+}\Gmat^\top\Wmat\,\bvec .
  \label{eq:slope}
\end{equation}
For $\bvec\in\rng(\Gmat)$ with $\Pi_{\mathrm{obs}}\Delta\hat{\xvec}\neq\bm0$ the denominator
vanishes ($\lVert\Pperp\bvec\rVert_{\Wmat}=0$ by Theorem~\ref{thm:undetectable}) while the
numerator does not, so $\operatorname{slope}(\bvec)=\infty$: the fault corrupts an observable
state combination while producing no parity residual. No finite alert limit protects this
fault class autonomously.
\end{proposition}

\begin{proof}
By Theorem~\ref{thm:undetectable}, $\bvec\in\rng(\Gmat)\Rightarrow\Pperp\bvec=\bm0\Rightarrow
\lVert\Pperp\bvec\rVert_{\Wmat}=0$. The estimator response
$\Delta\hat{\xvec}=(\Gmat^\top\Wmat\Gmat)^{+}\Gmat^\top\Wmat\bvec$ need not lie in
$\Nsp(\Gmat)$, so its observable projection $\Pi_{\mathrm{obs}}\Delta\hat{\xvec}$ can be
nonzero; the in-range witness of Section~\ref{sec:undetectable} (a node-1 offset shift, an
\emph{observable} relative-clock combination) is such a case. A ratio with zero denominator
and nonzero numerator is $+\infty$.
\end{proof}

The finite counterpart---the worst \emph{detectable} single-measurement slope and the
alert-limit crossing---is a \Modelled{} magnitude quantified in Section~\ref{sec:designlaws}.
